% =====================================================================
% Seção 4 — Estudo de Caso (R459/R460)
% =====================================================================
\section{Estudo de Caso: Diversificação em RAG}\label{sec:caso}

Apresentamos os resultados empíricos do benchmark R459/R460. Todos os valores
são do relatório auditável \texttt{cohort\_report.json}.

\subsection{Resultado agregado por estratégia}

A Tabela~\ref{tab:result} resume as métricas por estratégia, agregadas sobre as
4 consultas.

\begin{table}[h]
\centering
\small
\caption{Métricas agregadas por estratégia de diversificação no corpus-piloto
(4 consultas, top-$n$=8, $k$=4). Fonte: \texttt{cohort\_report.json}.}
\label{tab:result}
\begin{tabular}{lccc}
\toprule
\textbf{Estratégia} & \textbf{Div} & \textbf{Cob.} & \textbf{$g$} \\
\midrule
top-k (atual) & 0{,}5000 & 0{,}6667 & 0{,}4561 \\
MMR          & 0{,}7500 & 0{,}9167 & 0{,}4256 \\
Recamán      & 0{,}5000 & 0{,}6667 & 0{,}4561 \\
\textbf{HABD} & \textbf{0{,}8333} & \textbf{1{,}0000} & 0{,}4087 \\
\bottomrule
\end{tabular}
\end{table}

Três observações decorrem diretamente dos dados:

\begin{enumerate}
  \item \textbf{Recamán não diverge do top-k em diversidade} ($\Div = 0{,}5000$
        em ambos). Como o schema \emph{posicional} rearranja posições sem operar
        sobre âncoras, sob um ranking monopolista ele reproduz o top-k. O veredito
        é \texttt{refuta\_H2}: não há ganho estrito de diversidade sobre o
        top-k.
  \item \textbf{HABD maximiza diversidade e cobertura}: $\Div = 0{,}8333$ e
        $\cob = 1{,}0000$, superando MMR ($0{,}7500$ / $0{,}9167$) e Recamán. O
        veredito agregado registra \texttt{HABD supera MMR e Recamán em
        diversidade e cobertura}.
  \item \textbf{Mas a fidelidade cai além do tolerável}: a groundedness cai de
        $g_{\mathrm{base}} = 0{,}4561$ (top-k) para $g_{\mathrm{HABD}} =
        0{,}4087$, uma queda relativa de
        $\tfrac{\Delta g}{g} = 0{,}1039 > 0{,}05$. Formalmente:
        \[
        \frac{\Delta g}{g} \;=\;
        \frac{0{,}4561 - 0{,}4087}{0{,}4561} \;=\; 0{,}1039,
        \]
        que excede a tolerância de $5\%$ adotada.
\end{enumerate}

\subsection{Veredito das hipóteses}

Com base nos critérios definidos na Seção~\ref{sec:metodo}:

\textbf{H2 (Recamán supera o top-k em diversidade):}
\[
\begin{aligned}
\Delta\Div(\text{Recam\'an}) &= 0{,}0
  \quad \text{(n\~ao estritamente $>0$)}\\
&\Rightarrow\; \textbf{H2 refutada}
\end{aligned}
\]

\textbf{H3 (HABD supera sem sacrificar relevância):}
\[
\begin{gathered}
\Delta\Div(\mathrm{HABD}) > 0{,}5\\[2pt]
\Div(\mathrm{HABD}) \ge \Div(\mathrm{MMR}) \quad \text{(factual)}\\[2pt]
\tfrac{\Delta g}{g} = 0{,}104 > 0{,}05
   \;\Rightarrow\; \textbf{(cond. violada)}\\[2pt]
\Rightarrow \textbf{H3 refutada}
\end{gathered}
\]

Registramos \texttt{refuta\_H3}: o HABD \emph{supera} MMR e Recamán em
diversidade e cobertura, mas \emph{não satisfaz} o critério rígido de
relevância. As duas afirmações são simultaneamente verdadeiras e reportadas de
forma equilibrada: ganho observável na fronteira de diversidade, sem vitória
limpa no critério definido.

\subsection{Leitura como microcosmo do fator limítrofe}

O estudo de caso encarna, em escala menor, o fator limítrofe formalizado na
Seção~\ref{sec:metodo}. O HABD opera análogo a um ``gerador com âncora'': ele
\emph{não cria informação nova} do vazio, mas reorganiza itens a partir de
\emph{âncoras canônicas} já presentes no ranking (analogia a $\mathcal{D}_t^r$).
O ganho de diversidade é real, mas \emph{não é gratuito}: ele incide sobre a
relevância, exatamente como a teoria prevê que otimizar uma métrica à custa de
outra tem um custo de fidelidade. A métrica de \emph{groundedness} é a ``âncora
externa'' que expõe o custo --- sem ela, o ganho de diversidade seria reportado
como vitória limpa, o que constituiria um \emph{overclaim}.
